\section*{Results}

\subsection*{Most released architectures cannot forbid impossible tensor predictions}

We first audited which released architectures could, even in principle, guarantee a symmetry-mandated zero for a parity-odd tensor. Fifteen models were classified by reading their source at a pinned version or commit, recording the deciding line for each; where source inspection could not settle it, the model was probed numerically rather than guessed (Table~\ref{tab:audit}). Six attach parity labels to their features, so an odd-parity output vanishes structurally by construction. The nine that do not fail in three different ways, and the difference matters. Three are rotation-equivariant only: parity is absent from the features, so the guarantee cannot be recovered from them at all. Three carry a single unlabelled Cartesian vector channel, which is built from edge vectors and is therefore implicitly polar; a rank-3 head assembled from an odd number of vector factors would inherit the guarantee, but nothing in the framework types it or provides a rank-3 path, so that is the practitioner's to build and nothing checks it. Three are invariant models that cannot express an equivariant tensor output at all. Only the first group is beyond repair, and it is the group this paper measures. The rotation-only group is not a fringe. It contains EquiformerV2 and the eSCN family, on which the current generation of universal potentials is built\cite{fu2025esen,wood2025uma}. It also contains Equiformer v1, classed here by its released default configuration, whose node embeddings are all-even, although its formulation admits odd irreducible representations. The full audit, with the deciding source line for every row, is in Supplementary Note~2 and Supplementary Table~1.

Two rows show why measurement, not reading, has to settle such classifications. NequIP exposes a boolean named \texttt{parity}, yet its own documentation shows that disabling it does not remove parity labels from the features; the flag is not an SO(3) toggle, and our SO(3) arms are therefore built by relabelling irreducible representations directly. GotenNet, conversely, propagates features through attention blocks that are not parity-typed, and we expected them to discard parity. A reflection test showed the opposite: its vector output obeys the polar-vector law under improper operations to machine precision ($1.8\times10^{-15}$, against $1.9$ for the pseudovector law), so the row moved from undetermined to parity-aware on the measurement.

\begin{table}[t]
\centering
\caption{Feature parity across fifteen released architectures. Each model was classified by reading its released source at a pinned version or commit; two rows, ICTP and GotenNet, were settled instead by a numerical reflection test at random initialization, because the source did not decide them. Only the six architectures in the first group can make a parity-odd output vanish structurally on a centrosymmetric input. The three in the second group are the subject of this paper: their features contain no parity information at all, so the guarantee cannot be recovered from them. Versions, commits and the deciding source line for every row are given in Supplementary Table~1, and the classification does not extend to later releases.}
\label{tab:audit}
\fontsize{8bp}{10bp}\selectfont
\begin{tabular}{@{}l l c@{}}
\toprule
Architecture & Internal features & Parity typed \\
\midrule
\multicolumn{3}{@{}l}{\textbf{Parity-aware} (6). An odd-parity output is structurally zero.} \\
NequIP\cite{batzner2022nequip}      & O(3) irreps, natural parity $(-1)^l$ & yes \\
Allegro\cite{musaelian2023allegro}  & O(3) irreps, natural parity $(-1)^l$ & yes \\
MACE\cite{batatia2022mace}          & O(3) irreps, natural parity $(-1)^l$ & yes \\
MatTen\cite{wen2024matten}          & O(3) irreps, natural parity $(-1)^l$ & yes \\
ICTP\cite{zaverkin2024ictp}         & Cartesian harmonics; parity implicit & yes\textsuperscript{$\ast$} \\
GotenNet\cite{aykent2025gotennet}   & untyped attention; polar vector output & yes\textsuperscript{$\ast$} \\
\addlinespace[3pt]
\multicolumn{3}{@{}l}{\textbf{Rotation-only} (3). Nothing can force the output to vanish.} \\
Equiformer v1\cite{liao2023equiformer}   & SO(3) irreps, all declared even & no \\
EquiformerV2\cite{liao2024equiformerv2}  & degree and order only; no parity flag & no \\
eSCN\cite{passaro2023escn}               & degree and order only; no parity flag & no \\
\addlinespace[3pt]
\multicolumn{3}{@{}l}{\textbf{Vector-only} (3). Implicitly polar but untyped, and no rank-3 path.} \\
PaiNN\cite{schutt2021painn}          & one Cartesian vector channel & untyped \\
TorchMD-Net\cite{tholke2022torchmdnet} & one Cartesian vector channel & untyped \\
ViSNet\cite{wang2024visnet}          & one Cartesian vector channel & untyped \\
\addlinespace[3pt]
\multicolumn{3}{@{}l}{\textbf{Invariant} (3). Cannot express an equivariant tensor output.} \\
SchNet\cite{schutt2017schnet}            & interatomic distances only & n/a \\
DimeNet++\cite{gasteiger2020dimenetpp}   & distances and bond angles & n/a \\
FAENet\cite{duval2023faenet}             & not equivariant; frames applied externally & n/a \\
\bottomrule
\end{tabular}

\vspace{2pt}
{\fontsize{7bp}{8.5bp}\selectfont
\textsuperscript{$\ast$}Settled by measurement, not by reading the source. GotenNet's vector output obeys the polar-vector law under improper operations to $1.8\times10^{-15}$, against $1.9$ for the pseudovector law.\par}
\end{table}

\subsection*{What parity guarantees, and what rotation alone still guarantees}

Before any measurement, group theory fixes what each symmetry class can promise. We state the account in full here; the proofs are in Supplementary Note~1, and every statement is used by an experiment below.

\paragraph{Setting.} A periodic structure $x$ carries $n$ atoms at positions $\mathbf{r}_1,\dots,\mathbf{r}_n$ with species labels, together with a Bravais lattice $\Lambda$. A model $f_\theta$ with parameters $\theta$ maps $x$ into the target space $V$. For a rank-$r$ polar tensor, $V$ carries the representation $D(Q) = Q^{\otimes r}$ of $Q \in \mathrm{O}(3)$, so that $D(I) = (-1)^{r}\,\mathrm{id}$ at the inversion $I$, which sends every position to its negative. The piezoelectric tensor has $r = 3$ and $V = \mathbb{R}^{3}\otimes\mathrm{Sym}^{2}(\mathbb{R}^{3})$, the eighteen-dimensional space of rank-3 tensors symmetric in the strain indices, a symmetry that $D$ preserves. We write $\lVert\cdot\rVert_F$ for the Frobenius norm, so that $\lVert f_\theta(x)\rVert_F$ is the violation magnitude $\lVert e \rVert_F$ that every experiment below reports.

Three assumptions are used throughout. \textbf{(A1)} \emph{Permutation invariance}: relabelling atoms of the same species leaves the output unchanged. \textbf{(A2)} \emph{Translation invariance}: a rigid shift by any vector, and a shift of any atom by a lattice vector, leave it unchanged. \textbf{(A3)} \emph{$G$-equivariance}: $f_\theta(Q\!\cdot\!x)=D(Q)f_\theta(x)$ for all $Q\in G$, with $G=\mathrm{O}(3)$ for a parity-labelled network and $G=\mathrm{SO}(3)$ for one whose features carry no parity labels. The rigid-shift clause of (A2) is load-bearing, because a non-symmorphic space group pairs a point-group rotation with a fractional translation that is not a lattice vector. Every architecture studied here satisfies (A1) and (A2) by construction, through a permutation-invariant sum over per-atom or per-edge contributions on a lattice-periodic graph, and we verify both numerically (Supplementary Note~9).

\begin{definition}[Centrosymmetric structure]\label{def:centro}
$x$ is centrosymmetric if there is a point about which the inversion maps the structure onto itself. Placing that point at the origin, this means there is a species-preserving permutation $\pi$ of the atoms and lattice vectors $\mathbf{t}_i \in \Lambda$ with $-\mathbf{r}_i = \mathbf{r}_{\pi(i)} + \mathbf{t}_i$ for every $i$. Because $I^2 = \mathrm{id}$, the permutation is an involution, $\pi^{2} = \mathrm{id}$.
\end{definition}

One identity drives everything that follows, and it uses no assumption about parity.

\begin{lemma}[Inversion identity]\label{lem:inv}
Let $f_\theta$ satisfy (A1) and (A2). Then $f_\theta(I\!\cdot\!x) = f_\theta(x)$ for every centrosymmetric $x$ and every $\theta$.
\end{lemma}

The inverted copy of a centrosymmetric crystal is the same periodic solid, relabelled and shifted, so any model that cannot see relabelling or shifts cannot see the inversion either. By (A2) it does not matter about which point the inversion is taken: inverting about the cell origin, as the experiments below do, differs from inverting about the crystal's centre by a rigid shift. Adding parity to this identity gives the guarantee.

\begin{theorem}[Structural zero]\label{thm:zero}
Let $f_\theta$ satisfy (A1), (A2) and (A3) with $G = \mathrm{O}(3)$, and let the target be a rank-$r$ polar tensor with $r$ odd. Then $f_\theta(x) = 0$ for every centrosymmetric $x$ and every $\theta$.
\end{theorem}

The proof is two lines. Lemma~\ref{lem:inv} gives $f_\theta(I\!\cdot\!x)=f_\theta(x)$; (A3) at $Q=I$ gives $f_\theta(I\!\cdot\!x)=D(I)f_\theta(x)=-f_\theta(x)$. Hence $f_\theta(x)=-f_\theta(x)$, so $f_\theta(x)=0$. No step mentions the parameters, the data or the loss, which is what makes the zero structural rather than learned. For an SO(3)-labelled network the second step is simply unavailable: equivariance is asserted for proper rotations only, and $D(I)$ is never invoked. In exact arithmetic that is the whole story; in floating point, what is delivered instead is bounded.

\begin{proposition}[Arithmetic floor]\label{prop:floor}
Fix a centrosymmetric $x$. Let $\delta = \lVert f_\theta(I\!\cdot\!x) - f_\theta(x)\rVert_F$ measure the failure of (A1)--(A2) at $x$, and let $\eta = \lVert f_\theta(I\!\cdot\!x) - D(I) f_\theta(x)\rVert_F$ measure the failure of the O(3) law at $Q=I$, a law an SO(3)-labelled model never asserted. Then
$\lVert f_\theta(x)\rVert_F \le (\delta + \eta)/2$.
\end{proposition}

Writing $2f_\theta(x) = [f_\theta(x)-f_\theta(I\!\cdot\!x)] + [f_\theta(x)+f_\theta(I\!\cdot\!x)]$ and applying the triangle inequality proves it. Three things follow. With $\delta=\eta=0$ it recovers Theorem~\ref{thm:zero}. In floating point $\delta$ and $\eta$ are set by the arithmetic, so an O(3) model's residual floor tracks its precision and not its symmetry, which is why the floors sit at the float32 level and why MACE's is roughly eight times the others': its symmetric contractions stay in single precision even under a double-precision cast (Supplementary Note~3). And for a parity-blind model Lemma~\ref{lem:inv} forces $\delta = 0$ and $\eta = 2\lVert f_\theta(x)\rVert_F$, so the bound reads $\lVert f_\theta(x)\rVert_F \le \lVert f_\theta(x)\rVert_F$ and constrains nothing at all. The guarantee and its stability are lost together. Two consequences of Theorem~\ref{thm:zero} are tested directly below.

\begin{corollary}[Inversion averaging is vacuous]\label{cor:sym}
For any $f_\theta$ satisfying (A1)--(A2), parity-labelled or not, the antisymmetrized prediction $T_{\mathrm{sym}}(x) = [f_\theta(x) - f_\theta(I\!\cdot\!x)]/2$ vanishes identically on every centrosymmetric $x$.
\end{corollary}

This is Lemma~\ref{lem:inv} alone. The zero it produces carries no information about parity, holds whatever the model has learned, and degrades in proportion to $\delta$. The next extends the guarantee from the value to its derivative.

\begin{corollary}[The guarantee is differentiable]\label{cor:jac}
Under the hypotheses of Theorem~\ref{thm:zero}, with $f_\theta$ differentiable in the positions, let $P$ act on displacement fields $u\in\mathbb{R}^{3n}$ by $(Pu)_{j} = -u_{\pi^{-1}(j)}$, so that $P^{2}=\mathrm{id}$. The Jacobian $J = \partial f_\theta/\partial\mathbf{r}$ at a centrosymmetric $x$ satisfies $J\circ P = -J$, hence $J P_{\mathrm{even}} = 0$ with $P_{\mathrm{even}} = (\mathrm{id}+P)/2$: the response is supported entirely on inversion-odd displacement modes and the even-subspace energy fraction $\lVert J P_{\mathrm{even}}\rVert_F/\lVert J\rVert_F$ vanishes.
\end{corollary}

Finally, what does rotation alone still buy?

\begin{corollary}[Rotation ceiling]\label{cor:rot}
Let $f_\theta$ satisfy (A1)--(A2) and (A3) with $G=\mathrm{SO}(3)$ only, and let $G^{+}(x)\subset\mathrm{SO}(3)$ be the proper-rotation subgroup of the point group of $x$; each of its elements is the rotational part of a space-group operation of $x$, and so returns the crystal to itself once the accompanying translation, which may be fractional, is absorbed by (A2). Then $f_\theta(x)$ lies in the invariant subspace $V^{G^{+}(x)} = \{T \in V : R^{\otimes r}T = T\ \forall R \in G^{+}(x)\}$. If $V^{G^{+}(x)} = \{0\}$, then $f_\theta(x)=0$ is forced with no parity information; otherwise nothing forces it to vanish. Suppose instead that (A1)--(A2) and (A3) hold only approximately over $G^{+}(x)$, with defects
$\delta_{+} = \max_{R} \lVert f_\theta(R\!\cdot\!x) - f_\theta(x)\rVert_F$ and
$\epsilon = \max_{R} \lVert f_\theta(R\!\cdot\!x) - D(R) f_\theta(x)\rVert_F / \lVert f_\theta(x)\rVert_F < 1$.
Then the forced case weakens to $\lVert f_\theta(x)\rVert_F \le \delta_{+}/(1-\epsilon)$.
\end{corollary}

For the piezoelectric tensor this becomes a sharp, parameter-free prediction. Of the 32 crystal classes, 21 are non-centrosymmetric and 20 of those are piezoelectric; the one exception is class 432, whose rotations alone annihilate every rank-3 piezoelectric tensor\cite{nye1985crystals}. Exactly one of the eleven centrosymmetric classes, m$\bar 3$m, has 432 as its proper-rotation subgroup. Corollary~\ref{cor:rot} therefore forces a rotation-only model to zero on m$\bar 3$m crystals and on no others: 166 of our 2,000, leaving 1,834 (91.7\%) on which nothing but parity can enforce the zero. The stability clause explains why that is a ceiling only for an exactly equivariant model. The three matched SO(3) cores have $\delta_{+}$ and $\epsilon$ at the arithmetic floor and flag none of the 166 in any seed. EquiformerV2 has both of order $10^{-1}$, the bound is correspondingly weak, and it flags 53.4\% of them.

\paragraph{What the theory does not say.} It forbids; it does not compel. Theorem~\ref{thm:zero} says an O(3) model \emph{must} predict zero. Corollary~\ref{cor:rot} says an SO(3) model \emph{may} predict a nonzero tensor on 91.7\% of this population, not that it will. That it does, on 97.6 to 99.2\% of the crystals where it is free to, is an empirical finding, and it is what the rest of this paper measures.

\subsection*{A label-free test with matched-pair architectures}

The design is summarized in Fig.~\ref{fig:concept}. NequIP, Allegro and MACE were each built in two arms through a single irreducible-representation helper: an O(3) arm carrying the natural spherical-harmonic parity $(-1)^{l}$ on every feature, and an SO(3) arm in which every feature is declared parity-even. The geometric content of the two arms is identical; only the parity bookkeeping differs, and with it the reflection equivariance the feature algebra enforces. Architecture, hyperparameters, data, splits and seeds are shared. EquiformerV2 joins unmodified, because it has no parity labels to toggle, and that absence is the reason for its inclusion. The grid is 7 arms $\times$ 4 targets $\times$ 3 seeds $=$ 84 runs.

Before any training, every arm passed a verification gate probing an internal feature layer with a random rotation and a random improper operation: O(3) arms must satisfy both transformation laws, SO(3) arms must satisfy rotation and fail reflection. All arms behaved as required, and the SO(3) parameter count exceeds its O(3) counterpart in every configuration (ratios 1.00--1.40), so reduced capacity cannot explain any SO(3) deficit reported below (gate results in Supplementary Note~3 and Supplementary Table~2; per-configuration parameter counts in Supplementary Table~8).

Four targets separate two properties that a single experiment would confound: being a parity-odd target, and being evaluated where symmetry forbids a nonzero value. The QM9 internal energy $U_0$ (parity-even scalar) and the elasticity tensor (parity-even rank 4) carry no odd-parity constraint. The QM9 dipole is parity-odd, but no QM9 molecule has a point group that forces it to vanish, so the constraint is never active. Only the piezoelectric tensor on centrosymmetric crystals is both. The evaluation population is 2,000 centrosymmetric insulators from the Materials Project\cite{jain2013materials}, each re-verified with spglib\cite{togo2024spglib}, in two coordinate variants: idealized (snapped onto the exact space group) and raw (as relaxed by density functional theory).

\begin{figure}[t]
\centering
\includegraphics[width=\textwidth]{fig1_concept}
\caption{A label-free test of physical validity. \textbf{a}, A centrosymmetric crystal, shown with its unit cell (dashed) and two neighbouring cells. Every atom has a partner at $-\mathbf{r}$ of the same species, so the inversion maps the crystal onto itself, $I\!\cdot\!x = x$. The piezoelectric tensor is parity-odd, $e \mapsto -e$, and Neumann's principle requires it to be invariant under every operation of the point group; the two demands together force $e = -e$, hence $e = 0$ exactly. \textbf{b}, The two networks are built from one code path and differ in one bit. The O(3) arm labels its features with the natural spherical-harmonic parity $(-1)^{l}$, giving $0e,\,1o,\,2e,\,3o$ up to $l_{\max}=3$; the SO(3) arm declares every degree even. Architecture, hyperparameters, data, splits and seeds are identical, and the SO(3) arm has at least as many parameters (Supplementary Table~8). \textbf{c}, Predicted magnitudes on the 2,000 spglib-verified centrosymmetric crystals, for the NequIP pair; boxes span the 25th to 75th percentile, whiskers the 5th to 95th, and the horizontal line is the median. The true value is exactly zero for every structure. The O(3) arm returns the arithmetic floor of Proposition~\ref{prop:floor}, a median of $3\times10^{-7}$\,C\,m$^{-2}$, which is the finite-precision realization of an exact zero and not a small nonzero response. The SO(3) arm predicts tensors above the detection threshold, at the magnitude of a genuine piezoelectric.}
\label{fig:concept}
\end{figure}

\subsection*{Parity labels decide whether impossible predictions occur}

Figure~\ref{fig:headline}a shows the distribution of predicted piezoelectric magnitudes $\lVert e \rVert_F$ on the idealized evaluation population, where the true value is exactly zero for every structure. The two arms are separated by six orders of magnitude. Not one of the 18,000 O(3) structure-runs exceeds the operating threshold of 0.01\,C\,m$^{-2}$. Their medians sit at $3\times10^{-7}$ to $3\times10^{-6}$\,C\,m$^{-2}$, the float32 floor that Proposition~\ref{prop:floor} bounds by the arithmetic defects, and the largest single O(3) prediction anywhere is $1.3\times10^{-3}$\,C\,m$^{-2}$, a factor of seven below the threshold (Supplementary Note~4). The SO(3) arms exceed the same threshold for 89.5\% (NequIP), 91.0\% (Allegro), 90.8\% (MACE), and 95.7\% (EquiformerV2) of the population, which we call false flags, with median magnitudes of 0.44 to 0.96\,C\,m$^{-2}$. For calibration, 98.5\% of the real piezoelectric tensors in the training data exceed the threshold and their median magnitude is 0.51\,C\,m$^{-2}$: the typical SO(3) violation is the size of a genuine piezoelectric response, not a numerical residue.

The separation does not depend on the operating point. Over the two decades of threshold from $10^{-4}$ to $10^{-2}$\,C\,m$^{-2}$ the SO(3) false-flag fraction changes by at most 4 percentage points per core while the O(3) fraction never exceeds 0.0007 (Fig.~\ref{fig:headline}b); the full curves over 25 thresholds and both coordinate variants are given in Supplementary Note~4, Supplementary Fig.~1, and Supplementary Tables~3--4. Nor does it depend on the coordinate variant: the SO(3) false-flag fractions of the three matched cores are identical to four decimal places between idealized and raw coordinates, and EquiformerV2's changes by less than 0.001. The O(3) arms show a small nonzero raw-variant response traced to a single structure whose as-relaxed coordinates are genuinely non-centrosymmetric at tolerances below the selection tolerance, so the guarantee detected a data artefact rather than violating symmetry, as quantified in Supplementary Note~5 and Supplementary Table~5.

Paired by structure, one-sided Wilcoxon signed-rank tests on the 2,000 paired violations return $p$-values below the double-precision underflow limit ($p<10^{-300}$) for every core and seed, with the O(3) violation smaller for 96--98\% of individual crystals (Supplementary Note~6 and Supplementary Table~6). The violation pattern is also not an idiosyncrasy of any one model. Per-structure violations correlate between independently seeded runs of the same architecture (Spearman $\rho = 0.77$--$0.83$), and independent architectures flag nearly the same crystals: the flagged sets of any two SO(3) cores overlap with Jaccard index 0.93--0.99. The failure is a property of the symmetry class of the model, not of a particular implementation.

\begin{figure}[t]
\centering
\includegraphics[width=\textwidth]{fig2_headline}
\caption{Parity labels separate possible from impossible predictions by six orders of magnitude. \textbf{a}, Distribution of the predicted piezoelectric magnitude $\lVert e \rVert_F$ on 2,000 spglib-verified centrosymmetric crystals (idealized coordinates), for the O(3) and SO(3) arms of each core. The true value is exactly zero for every structure, so the shaded region above the operating threshold $\tau$ is physically impossible. Boxes span the 25th to 75th percentile of the per-structure distribution (mean of per-seed percentiles over 3 seeds), whiskers the 5th to 95th, the horizontal line is the median, and the isolated dash is the maximum. Percentages above each SO(3) box give the fraction of structures exceeding $\tau$. The double arrow gives the ratio of the two NequIP medians, $2\times10^{6}$. \textbf{b}, False-flag fraction as a function of threshold (idealized coordinates), mean over 3 seeds; shaded bands show one standard deviation. Blue, O(3) arms; orange, SO(3) arms; markers distinguish cores. EquiformerV2 has no O(3) arm.}
\label{fig:headline}
\end{figure}

\subsection*{The constraint costs nothing where it is inactive}

If parity labelling handicapped models on ordinary tasks, removing it could be a defensible trade. It does not. On the three targets where no symmetry-forbidden value is ever evaluated, namely $U_0$, dipole, and elasticity, the two arms are indistinguishable within seed noise: the between-arm difference in test mean absolute error (MAE) is at most 0.94 pooled seed standard deviations, against a seed spread that is itself of that order (Fig.~\ref{fig:accuracy}a; full accuracy tables in Supplementary Note~7 and Supplementary Table~7). The differences are small but not sign-random. On $U_0$, the target with by far the largest seed variance, the sign varies across cores ($-0.42$ to $+0.94$ pooled standard deviations). On the dipole and on elasticity, the O(3) arm is nominally the better of the two in all three cores, by 0.04 to 0.74 pooled standard deviations. Three seeds cannot resolve an effect of that size, and we claim none; we note only that where any tendency is visible it favours keeping the parity labels, never removing them.

The dipole is the decisive control. It is parity-odd, like the piezoelectric tensor, yet the arms are within seed noise of each other, because QM9 contains no molecule whose symmetry forces the dipole to vanish: the constraint an O(3) model enforces is never active there. Of the two properties the piezoelectric experiment would otherwise confound, an odd-parity target and an evaluation where symmetry forbids a nonzero value, only the second produces a difference between the arms.

Where the constraint is active, it helps rather than costs. On the piezoelectric tensor the O(3) arm has lower in-distribution test error than the SO(3) arm for all three matched cores, by 3.4 to 4.6 pooled seed standard deviations with a consistent sign (Fig.~\ref{fig:accuracy}b). EquiformerV2's in-distribution error is comparable to the matched cores, which matters for interpretation: its near-universal false flagging on centrosymmetric crystals coexists with ordinary regression competence on non-centrosymmetric ones.

\begin{figure}[t]
\centering
\includegraphics[width=\textwidth]{fig3_accuracy}
\caption{Removing parity labels costs nothing where the parity constraint is inactive, and costs accuracy where it is active. \textbf{a}, Difference in test mean absolute error (MAE) between the SO(3) and O(3) arms, in units of the pooled per-seed standard deviation, for the three matched cores on all four targets. Positive values favour the O(3) arm. The grey band marks one pooled standard deviation, within which a three-seed comparison cannot resolve a difference; the tinted column is the one target on which the parity constraint is active. On the dipole and on elasticity the O(3) arm is nominally better in all three cores (3/3), by 0.04 to 0.74 pooled standard deviations, an effect three seeds cannot resolve and from which we claim no significance; on $U_0$, the target with by far the largest seed variance, the sign varies across cores. \textbf{b}, Piezoelectric test MAE for both arms of each core (mean over 3 seeds; error bars, one standard deviation), with the two arms of a core joined and the effect size printed above. EquiformerV2 has no O(3) arm. MAE, mean absolute error.}
\label{fig:accuracy}
\end{figure}

\subsection*{Training cannot buy the zero}

The natural objection is that the SO(3) failure is a data problem: no centrosymmetric crystal appears in the piezoelectric training set, so the model was never told the answer. We tested this directly by retraining every SO(3) arm on the training set augmented with 1,000 fresh spglib-verified centrosymmetric insulators labelled with the exact zero tensor, drawn from six centrosymmetric space groups, with hyperparameters, validation, and test partitions unchanged.

The decisive measurement is the in-training control: predictions on the very crystals the model was trained, with explicit zero labels, to call zero (Fig.~\ref{fig:training}a). After training, the models still false-flag 89.5--92.2\% of those crystals. The failure is not underfitting, since the zero-labelled rows are fit better than the real-tensor rows in training MAE, and it is not a data-coverage problem, since held-out false-flag fractions barely move whether or not the held-out crystal's space group appeared in training (changes of 0.00--0.03; full tables in Supplementary Note~8 and Supplementary Tables~10--11). Gradient descent drives the predictions towards zero, halving the violation medians, and cannot arrive, because the required value is exact. Augmentation meanwhile improves the ordinary regression error on non-centrosymmetric test crystals for every core: it produces better regressors that still predict impossible values.

Weighting the zero-labelled rows more heavily in the loss does not change the conclusion (Fig.~\ref{fig:training}b). Up-weighting them one-hundred-fold reduces the false-flag fraction on the trained-on crystals only from 0.90 to 0.68, leaves 87\% of held-out crystals from the same space groups flagged, and never approaches the exact zero the O(3) arms deliver at any weight, for free: the violation median on those crystals falls from 0.120 to 0.014\,C\,m$^{-2}$, still more than four orders of magnitude above the O(3) floor (Supplementary Table~12).

A second proposed repair operates at test time: antisymmetrize the output over inversion, $T_{\mathrm{sym}}(x) = [T(x) - T(I\!\cdot\!x)]/2$. On the exactly-equivariant SO(3) cores this drives the false-flag fraction to zero, and Corollary~\ref{cor:sym} says why the success is empty: by Lemma~\ref{lem:inv} the construction vanishes identically for \emph{any} permutation- and translation-invariant model, whatever it has learned, so no parity information is involved. The construction also presupposes knowing the target's parity in advance, which is precisely the knowledge O(3) features carry, and it fails on EquiformerV2, whose approximate equivariance violates the permutation identity the fix relies on: its false-flag fraction falls only to 0.82 (Supplementary Note~9 and Supplementary Table~13).

\begin{figure}[t]
\centering
\includegraphics[width=\textwidth]{fig4_training}
\caption{Zero-labelled training data and loss reweighting reduce, but cannot eliminate, impossible predictions. \textbf{a}, False-flag fraction on the 1,000 centrosymmetric crystals that each SO(3) arm was trained on with explicit zero labels, after augmented retraining (mean over 3 seeds). Every arm still flags roughly nine crystals in ten of those it was shown the answer for. The blue line marks the O(3) arms, which predict zero structurally and were not retrained. \textbf{b}, False-flag fraction of the NequIP SO(3) arm as the loss weight on the zero-labelled rows increases from 1 to 100, evaluated on the trained-on crystals and on held-out centrosymmetric crystals from space groups seen or unseen during training (mean over 3 seeds; error bars, one standard deviation). A hundredfold weight buys a reduction of 0.22 on the trained-on crystals and never approaches the exact zero that the O(3) arms deliver at any weight, for free.}
\label{fig:training}
\end{figure}

\subsection*{The guarantee tracks the physics of symmetry breaking}

Three measurements test the theory directly.

First, the guarantee switches off when the symmetry does. Displacing rutile TiO$_2$ along a polar mode, over 33 amplitudes from the centrosymmetric parent to beyond the physical distortion and with every frame symmetry-verified, the O(3) arms rise smoothly from the arithmetic floor of Proposition~\ref{prop:floor} ($1.9\times10^{-7}$ to $9.1\times10^{-7}$\,C\,m$^{-2}$ at zero distortion), while the SO(3) arms sit on a spurious floor of 0.09--0.28\,C\,m$^{-2}$ that is present already at zero distortion and dominates the physical signal throughout the small-distortion regime (Supplementary Fig.~2). Rutile carries the argument because its rotation subgroup 422 is piezoelectric, so by Corollary~\ref{cor:rot} only parity forbids a response there. \MISSING{Supplementary Fig.~2 cannot be drawn: the 33-amplitude sweep array is not in the manuscript. Export it to figdata/fig5a\_rutile\_sweep.csv and rerun build\_figures.py. The numbers quoted in this paragraph are unaffected}

Second, Corollary~\ref{cor:jac} requires the even-subspace energy fraction of the Jacobian to be exactly zero at a centrosymmetric geometry. Measured over 20 crystals spanning 19 space groups and all seeds, with every Jacobian certified against finite differences, it is $1.4\times10^{-7}$ to $9.7\times10^{-7}$ for the O(3) arms and 0.42--0.54 for the SO(3) arms: five to six orders of magnitude apart, with no overlap over any structure or seed (Fig.~\ref{fig:mechanism}b). The O(3) architecture does not merely output a zero; it encodes which atomic displacements can activate a piezoelectric response.

Third, the crystals SO(3) models get right are exactly the ones Corollary~\ref{cor:rot} identifies (Fig.~\ref{fig:mechanism}a). The three matched cores false-flag none of the 166 m$\bar 3$m crystals in any seed, as the corollary requires, and flag 89--100\% of the 18 crystals of class m$\bar 3$, whose rotation subgroup 23 is piezoelectric (per-seed range; seed means 94.4--98.2\%). Both rank statements are confirmed by Reynolds-operator tests in the released code. Across the population they flag 97.6\%, 99.2\% and 99.0\% of the 1,834 crystals on which nothing forces them to zero: an exactly rotation-equivariant model saturates the ceiling and violates the constraint essentially everywhere it is free to. EquiformerV2 exceeds it, 95.7\% against 91.7\%, as the stability clause allows: it flags 53.4\% of the m$\bar 3$m crystals an exact model cannot flag at all, worth 4.4 percentage points, the entire excess. SO(3) models keep the zeros rotations already guarantee and only those (Supplementary Note~10 and Supplementary Table~15).

\MISSING{the state-of-the-art crystal-tensor predictors that enforce symmetry at the readout (GMTNet, GoeCTP, EATGNN, CEITNet) are criticized in the Discussion and never tested. Run all four on the same 2,000 centrosymmetric crystals and report their false-flag fractions here, alongside the even-subspace energy fraction of their Jacobians. If they output exactly zero, the field already solves the output problem and the contribution must be renarrowed, honestly, to the features underneath. If they output nonzero, the headline extends to the dedicated symmetry-enforcing state of the art, which is a much larger result. Either outcome strengthens the paper; the current omission leaves it exposed to the first referee drawn from that community. See EXPERIMENTS\_TO\_RUN.md, B-1}

The pattern is the same on familiar named materials (Fig.~\ref{fig:mechanism}c): the SO(3) arms predict 0.21 to 1.13\,C\,m$^{-2}$ for corundum Al$_2$O$_3$ (sapphire), whose non-cubic point group leaves parity as the only protection, and near zero for the rocksalt and diamond-structure compounds of class m$\bar 3$m (Supplementary Table~14).

\begin{figure}[t]
\centering
\includegraphics[width=\textwidth]{fig5_mechanism}
\caption{The missing guarantee is specifically parity, and it is differentiable. \textbf{a}, False-flag fraction of the SO(3) arms split by the point-group family of the evaluated crystal: m$\bar 3$m, whose proper-rotation subgroup 432 alone forbids a rank-3 piezoelectric tensor; m$\bar 3$, whose subgroup 23 does not; and all non-cubic classes, none of whose rotation subgroups forbid it. The three matched cores flag none of the 166 m$\bar 3$m crystals in any seed, as Corollary~\ref{cor:rot} requires. EquiformerV2's nonzero bar there reflects its approximate rotational equivariance. The O(3) arms are zero in every family. \textbf{b}, Even-subspace energy fraction of the Jacobian $J = \partial T/\partial \mathbf{r}$ at centrosymmetric geometries of 20 crystals spanning 19 space groups, measured over both arms of the three matched cores and all seeds; the boxes give the full measured range, which Corollary~\ref{cor:jac} forces to zero for an O(3) model. The two arms are five to six orders of magnitude apart with no overlap over any structure or seed. \textbf{c}, Predicted magnitude on ten named compounds (Supplementary Table~14). Shaded rows are class m$\bar 3$m, where rotation alone already forbids a response; on the two unshaded, non-cubic rows only parity protects the zero, and every SO(3) arm fails there. On the shaded rows, rotation alone already forbids a response, so the three exactly rotation-equivariant SO(3) cores return a true zero while the O(3) markers sit at their arithmetic floor of order $10^{-7}$\,C\,m$^{-2}$. EquiformerV2, whose rotational equivariance is only approximate, nonetheless exceeds the threshold on two of them, diamond and SrTiO$_3$, which is the same failure its 53.4\% false-flag rate on m$\bar 3$m records in \textbf{a}. The two unshaded rows are the ones that discriminate parity from rotation. Filled circle, O(3); open markers, the SO(3) arms. The axis is symmetric-log, so an exact zero is plotted at zero rather than clamped to a small positive value. \VERIFY{Supplementary Table~14 prints the SO(3) entries on the m$\bar 3$m rows to three decimals as 0.000, and this panel places them at an exact zero on the strength of Corollary~\ref{cor:rot} and Supplementary Table~15, which reports 0.0000 for that family. Confirm from the run outputs that they are exact zeros and not values rounded down from below $5\times10^{-4}$}}
\label{fig:mechanism}
\end{figure}
